\section{Architecture Integration}

\subsection{Model Support}

\begin{table}[H]
\centering
\small
\begin{tabular}{llllll}
\toprule
\textbf{Model} & \textbf{Attention} & \textbf{Nested?} & \textbf{RoPE?} & \textbf{GQA?} & \textbf{Conv1D?} \\
\midrule
DistilBERT & \texttt{MultiHeadSelfAttention} & No & No & No & No \\
RoBERTa & \texttt{BertSelfAttention} & Yes & No & No & No \\
GPT-2 & \texttt{GPT2Attention} & No & No & No & Yes \\
Qwen3.5 & \texttt{Qwen3Attention} & No & Yes & Yes & No \\
\bottomrule
\end{tabular}
\end{table}

\subsection{Weight Transfer Details}

\subsubsection{GPT-2 Conv1D}

\begin{lstlisting}
if hasattr(src, "c\_attn"):
    weight = src.c\_attn.weight.data
    if weight.shape[1] == 3 * self.hidden\_size:
        weight = weight.t()
    q\_w, k\_w, v\_w = weight.chunk(3, dim=0)
\end{lstlisting}

The Conv1D stores weights as \texttt{(in, out)}, opposite to \texttt{nn.Linear}. We 
transpose before chunking.

\subsubsection{RoBERTa Nested}

\begin{lstlisting}
if hasattr(src, "self"):
    s = src.self
    extracted["q"] = (s.query.weight, ...)
    extracted["k"] = (s.key.weight, ...)
    extracted["v"] = (s.value.weight, ...)

if hasattr(src, "output") and hasattr(src.output, "dense"):
    extracted["o"] = (src.output.dense.weight, ...)
\end{lstlisting}

\subsubsection{Qwen GQA}

\begin{lstlisting}
named\_map = {
    "q": ["q\_proj"],
    "k": ["k\_proj"],  # (H\_kv d\_h, d)
    "v": ["v\_proj"],
    "o": ["o\_proj"],
}
\end{lstlisting}

The K, V projections are smaller in GQA ($H_{kv} < H$).

\subsubsection{DistilBERT}

\begin{lstlisting}
named\_map = {
    "q": ["q\_lin"],
    "k": ["k\_lin"],
    "v": ["v\_lin"],
    "o": ["out\_lin"],
}
\end{lstlisting}

\subsection{Nested Encoder Bypass}

For nested encoders:

\begin{lstlisting}
if self.is\_nested\_encoder:
    output = context  # Outer layer applies output projection
else:
    output = self.out\_proj(context)
\end{lstlisting}

\subsection{LoRA on Nested Output}

For nested encoders, \texttt{out\_proj} is set to identity:

\begin{lstlisting}
elif proj\_key == "o" and self.is\_nested\_encoder:
    min\_d = min(target.weight.shape[0], target.weight.shape[1])
    target.weight.data.zero\_()
    target.weight.data[:min\_d, :min\_d].fill\_diagonal\_(1.0)
    if target.bias is not None:
        target.bias.data.zero\_()
\end{lstlisting}

This ensures the LoRA branch on \texttt{out\_proj} starts from identity, and the 
outer layer's projection dominates.

\subsection{Injection Algorithm}

\textbf{Algorithm 2:} Inject LoRA Attention
\begin{lstlisting}
Input:  model M, r\_L, alpha
Output: M with attention replaced by LoRA modules

1: for   in transformer\_layers(M) do
2:     src $\leftarrow$ get\_attention( )
3:     new $\leftarrow$ LoRAAttention(cfg, r\_L, alpha, pretrained\_attn=src)
4:     replace( , new)
5: end for
6: return M
\end{lstlisting}

\subsection{Failure Modes and Fallbacks}

\begin{table}[H]
\centering
\small
\begin{tabular}{lll}
\toprule
\textbf{Failure} & \textbf{Detection} & \textbf{Response} \\
\midrule
Missing Q/K/V & \texttt{missing} non-empty & \textbf{Raise RuntimeError} \\
Near-zero norm & Frobenius check & \textbf{Raise RuntimeError} \\
copy\_weights fails & Exception & \textbf{Raise RuntimeError} \\
Injection fail & Exception & \textbf{Raise RuntimeError} \\
\bottomrule
\end{tabular}
\end{table}

\textbf{Note.} Unlike other modules in this suite, LoRA Attention **does not 
gracefully degrade\textbf{ -- it fails hard. This is intentional because }silent 
failure would corrupt the base weights**.

\subsection{Registry}

\begin{lstlisting}
@ATTENTION\_REGISTRY.register("lora\_attn")
@ATTENTION\_REGISTRY.register("lora\_attention")
class LoRAAttention(BaseAttention):
    ...
\end{lstlisting}

\subsection{Interaction with Compression}

LoRA Attention can be applied on top of any compressed attention:

\begin{lstlisting}
# 1. Compress with SVD
compressed\_attn = SVDLinear(original\_attn, rank=64)

# 2. Wrap with LoRA
lora\_attn = LoRALinear(compressed\_attn)

# 3. Fine-tune LoRA only
for p in compressed\_attn.parameters():
    p.requires\_grad = False
for p in lora\_attn.parameters():
    p.requires\_grad = True
\end{lstlisting}

This pattern enables \textbf{post-compression recovery} for SVD/pruning/quantization 
methods.

\subsection{Merge Procedure}

\begin{lstlisting}
def merge\_lora(model):
    """Merge LoRA adapters into base weights for zero inference overhead."""
    for name, module in model.named\_modules():
        if isinstance(module, LoRALinear):
            with torch.no\_grad():
                delta = (module.lora\_B @ module.lora\_A) * module.lora\_scaling
                module.weight.data += delta
                # Delete LoRA params
                module.lora\_A = None
                module.lora\_B = None
                module.lora\_rank = 0
\end{lstlisting}

After merging, \texttt{model} has no LoRA parameters, only the merged weights.
\begin{lstlisting}

---
\end{lstlisting}

\begin{figure}[H]
\centering
\includegraphics[width=\linewidth]{figures/architecture.png}
\caption{High-level architectural diagram of lora\_attention.}
\label{fig:arch\_lora-attention}
\end{figure}

